\subsection{Installation des prérequis}

\begin{tip}
    Le serveur est supposé installé avec Alma 9 en configuration minimale.
\end{tip}

\subsubsection{Configuration du disque supplémentaire}
Le disque supplémentaire est celui contenant les fichiers web.\\

\begin{itemize}
    \item La première ligne spécifie le disque à converstir en volume physique (PV) ; ne pas oublier de remplacer la valeur présentée par le véritable nom du disque
    \item La seconde ligne crée un groupe de volumes (VG) nommé \texttt{VG\_Repo} à partir du volume physique créé précédemment
    \item La troisième ligne crée un volume logique (LV) nommé \texttt{/dev/VG\_Repo/LV\_repo} occupant tout l'espace libre du groupe de volumes
    \item La quatrième ligne formate le volume logique en système de fichiers XFS
    \item La cinquième ligne crée le point de montage \texttt{/var/www} pour le volume logique
    \item La sixième ligne ajoute une entrée dans le fichier \texttt{/etc/fstab} pour monter automatiquement le volume logique au démarrage
    \item La septième ligne monte tous les systèmes de fichiers définis dans \texttt{/etc/fstab}
    \item La dernière ligne recharge la configuration du démon systemd pour prendre en compte les modifications apportées au fichier \texttt{/etc/fstab}
\end{itemize}

\begin{lstlisting}[language=bash, caption=Configuration du disque supplémentaire]
pvcreate /dev/nvme0n2
vgcreate VG_Repo /dev/nvme0n2
lvcreate -n /dev/VG_Repo/LV_repo -l +100%FREE
mkfs -t xfs /dev/mapper/VG_Repo-LV_repo
mkdir /var/www
echo -e "/dev/mapper/VG_Repo-LV_repo               /var/www       xfs  nosuid,noexec,nodev        0 0" >> /etc/fstab
mount -a
systemctl daemon-reload
\end{lstlisting}

\subsubsection{Installation des prérequis applicatifs}
Les disques installés, les prérequis doivent désormais être installés.\\

\begin{itemize}
    \item Installer Apache, le support SSL, les outils de synchronisation, de téléchargement et les outils DNF
    \item Supprimer la page de démarrage d'Apache (lignes 2 à 14)
    \item Activer et démarrer le service Apache
\end{itemize}

\begin{lstlisting}[language=bash, caption=Installation des prérequis applicatifs]
dnf -y install httpd mod_ssl dnf-utils wget rsync
echo "# This file is part of Infr@Home Information System
# It is under GPL-V3 license
#
# Versionning
# YYYYMMDD_hhmm | Author                | Changelog
# 20240525_0957 | jmy37                 | File creation
#
# This configuration file enables the default "Welcome" page if there
# is no default index page present for the root URL.  To disable the
# Welcome page, comment out all the lines below.
#
# NOTE: if this file is removed, it will be restored on upgrades.
#" > /etc/httpd/conf.d/welcome.conf
systemctl enable --now httpd
\end{lstlisting}

\subsubsection{Configuration du pare-feu}

\begin{important}
Convertir cette syntaxe en une syntaxe iptables.
\end{important}

Les flux requis doivent être ouverts sur le pare-feu local.

\begin{lstlisting}[language=bash, caption=Configuration du pare-feu local]
firewall-cmd --add-service={ssh,http,https} --permanent
firewall-cmd --reload
\end{lstlisting}